% ==============
\chapter{Motion}
% ==============
%\section{Bodies and Particles}

%Let body $\cB$ be an open, bounded subset of three-dimensional           
%Euclidean space $\realnsd$ with boundary $\partial \cB$.                 
%The number of space dimensions, $\nsd$, is assumed equal                    
%to three unless stated otherwise.
%The boundary is decomposed within a particular space dimension into         
%two non-overlapping subdomains, $\Gamma_u$ and $\Gamma_{\sigma}$        
%where displacements or tractions are prescribed.  Thus                      
%$\partial \cB = \Gamma_u \union \Gamma_{\sigma}$ and 
%$\Gamma_u \intersection \Gamma_{\sigma} = \emptyset$.           
%Set closure is denoted                                                      
%$\bar{\cB} = \cB \union \partial \cB$.                             
%Let the body be composed of a closed set of particles with 
%coordinates $\vX$ in the reference configuration.  The particle
%has coordinates $\vX = \{X_1, X_2, X_3\}^T$ as measured from
%some origin $O$ in the dextral, orthonormal inertial reference
%%from $\veEihat, \veEjhat, \veEkhat$ attached to inertial
%from $\vefa, \vefb, \vefc$ attached to inertial
%reference frame $\cF$.


\section{General Motion}
Let the arbitrary time interval be defined as
$\interval \defe [t_0,t_f]$, initial to final
time, inclusive.\footnote{Note that $t_0$, while
typically zero, may be any real number less than $t_f$.}  
Let the {\bf motion}, a one parameter family of configurations, 
$\vvarphi : \cB \times \interval \mapsto \realthree$ 
map the material particle
$\vX$ into the current configuration $\vx$ such that
\be
\vx = \vvarphi(\vX,t) = \vvarphi_t(\vX).
\label{eq:motion}
\ee

A motion $\vvarphi$ evaluated
at a particular time $t \in \interval$ 
is referred to as a {\bf configuration} or
{\bf placement}.  To each configuration, we assign a {\bf displacement}
field $\vu: \cB \mapsto \realthree$ such that 
\be
\vu(\vX,t) \defe \vvarphi(\vX,t) - \vX 
\;\;\Longleftrightarrow\;\;
u_i \defe \varphi_i - \delta_{iJ} X_J.
\label{eq:displacement}
\ee
Thus, the current configuration $\vvarphi$ is simply a function 
of the original
placement $\vX$, 
plus a displacement $\vu$, which is {\bf function}
of placement $\vX$ and time $t$:
\be
\vvarphi(\vX,t) = \vX + \vu(\vX,t)
\;\; \Longleftrightarrow \;\;
\varphi_i = \delta_{iJ} X_J + u_i.
\label{eq:displacement_vvarphi}
\ee
It is important to
emphasize that the displacement remains a function $\vu = \vu(\vX, t)$
of reference configuration $\vX$ and time $t$.  
%
% This is call {\bf translation}.  
%A specialized case of translation, 
%where the displacement
%function $\vu = \bar{\vu}$ is constant is called {\bf offset}.
%See Section~\ref{section:offset} for more details.


\section{Deformation Gradient}
To each configuration $\vvarphi$, we define a deformation gradient 
$\tFbold: \vB \times \interval \mapsto \Mthreeplus$ 
such that
\begin{equation}
\tFbold 
\defe 
\Grad \vvarphi(\vX, t)
\;\; \Longleftrightarrow \;\;
\tF_{iJ} 
\defe 
\frac{\partial \varphi_i}{\partial X_J}.
\label{eq:deformation_gradient}
\end{equation}
%Remarks:
\begin{Hremark}
\begin{itemize}
\item Lower case indices represent components of tensors in the
current configuration whereas upper case indices represent
components of tensors in the reference configuration.  
\item Real, square matrices of dimension three and positive determinants
are denoted $\Mthreeplus$.  
\item $\Grad(\cdot) = \partial (\cdot) / \partial X_J$ denote
gradients taken in the reference configuration.
\item $\grad(\cdot) = \partial (\cdot) / \partial x_j$ denote
gradients taken in the current configuration. 
\end{itemize}
\end{Hremark}

From the displacement field in 
(\ref{eq:displacement}), we see that 
the relationship between the displacement gradient and the 
deformation gradient is given by
\be
\Grad \vu = \Grad \vvarphi - \vid = \tFbold - \vid
\;\; \Longleftrightarrow \;\;
\dst\frac{\partial u_i}{\partial X_J} = 
\dst\frac{\partial \varphi_i}{\partial X_J} - 
\delta_{iJ}.
\label{eq:Gradu}
\ee

\section{Jacobian of the Deformation Gradient}
The Jacobian of the deformation gradient, 
\begin{equation}
 J \defe \det \tFbold,
\end{equation}
describes the (generally non-uniform) volumetric expansion or 
contraction of the motion from the reference configuration $\vX$.

%\noindent Remarks:
\begin{Hremark}
\begin{itemize}
\item All configurations must be admissible in the sense that the Jacobian
of the deformation must be positive ($J > 0$).
This requirement keeps the deformations from mapping the body to
a single, infinitesimally small point ($J=0$) or turning the
body inside-out ($J<0$).  
\item Isochoric motions preserve volume, thus a Jacobian
of unity ($J=1$) describes an isochoric motion.
\item Table~\ref{tab:jacobian_categories} describes the categories
of motions (deformations) by Jacobian measure.  
\end{itemize}
\end{Hremark}

\begin{table}[htb]
\caption{Jacobian measure to categorize deformations.}
\label{tab:jacobian_categories}
\begin{center}
\begin{tabular}{|c|c|c|c|c|} \hline
 $J<0$ & $J=0$ & $0<J<1$ & $J=1$ & $J>1$ \\ %\hline \hline
 {\bf inadmissible}
 & {\bf inadmissible}
 & {\bf contraction} 
 & {\bf isochoric} 
 & {\bf expansion} 
 \\ \hline \hline
  &  &  & translation &  \\ 
  &  &  & rotation &  \\ 
  &  &  & isochoric stretch & \\ 
  &  &  & isochoric shear   & \\ \hline
 body 
 & body 
 & body's  
 & body's 
 & body's
 \\
 has turned 
 & has shrunk 
 & total volume 
 & total volume
 & total volume \\ 
 inside-out 
 & to zero volume 
 & has decreased 
 & is preserved 
 & has increased \\ \hline
\end{tabular}
\end{center}
\end{table}




